% Generated draft: requires mathematical and source review.
\documentclass[11pt]{article}
\usepackage[margin=25mm]{geometry}
\usepackage{fontspec}
\setmainfont{Latin Modern Roman}
\usepackage{amsmath,amssymb,mathtools,unicode-math}
\setmathfont{Latin Modern Math}
\usepackage{enumitem,needspace,xurl,hyperref}
\setlength{\parindent}{0pt}
\setlength{\parskip}{5pt}
\setlength{\emergencystretch}{3em}
\newcommand{\N}{\mathbb{N}}
\newcommand{\Z}{\mathbb{Z}}
\newcommand{\Q}{\mathbb{Q}}
\newcommand{\R}{\mathbb{R}}
\newcommand{\C}{\mathbb{C}}
\newcommand{\F}{\mathbb{F}}
\newcommand{\A}{\mathbb{A}}
\newcommand{\PP}{\mathbb{P}}
\newcommand{\E}{\mathbb{E}}
\newcommand{\eps}{\varepsilon}
\newcommand{\dd}{\,\mathrm{d}}
\begin{document}
\section*{UnsolvedMath ID: 10600012}
\textbf{Source problem code:} catalogue entry
\subsection{Definitions and mathematical statement}
An ordinary quandle has an idempotent, right-invertible, right-self-distributive operation \(*\). In one common convention, a virtual quandle is a quandle \(Q\) equipped with an automorphism \(f:Q\to Q\), so that \(f\) is bijective and
\[
 f(x*y)=f(x)*f(y).
\]
This convention is explanatory; the supplied record does not define its intended virtual quandle or virtual-crossing action.

Ordinary quandle homology is obtained from the rack tuple complex by quotienting out tuples with two adjacent equal entries. In a standard integral convention, its boundary before that quotient is
\[
 \partial(x_1,\ldots,x_n)=\sum_{i=2}^{n}(-1)^i
 \bigl((x_1,\ldots,\widehat{x_i},\ldots,x_n)
 -(x_1*x_i,\ldots,x_{i-1}*x_i,x_{i+1},\ldots,x_n)\bigr).
\]
Coefficient and low-degree conventions must be fixed separately.

Problem-list item 12 asks to study virtual quandle homology in analogy with quandle homology. It supplies no chain groups, differential, degeneracy quotient, coefficient category, or proposed theorem for the virtual version. It also gives no specified comparison with ordinary quandle homology or invariant to recover. Thus this is a development program whose defining construction remains missing, not a request to compute a source-defined complex already given here. In particular one cannot simply choose to ignore \(f\), take an invariant subcomplex, or form a mapping cone without further source evidence.
\subsection{Short English statement}
Develop virtual-quandle homology by analogy with ordinary quandle homology. The record does not define the virtual object or homological complex, so no particular extension is selected. Its supplied no-dedicated-literature assessment is a limited search report, not a verified absence claim.
\subsection{Sources}
\begin{itemize}
\item Local UnsolvedMath registry, permanent ID \(10600012\), exact source code \texttt{catalogue entry}, ``Virtual-knot problem 12 --- Virtual Quandle Homology.'' No individual proposer or proposed year is supplied.
\item \emph{Unsolved Problems in Virtual Knot Theory and Combinatorial Knot Theory}, revised survey, dedicated problem-list item 12: \url{https://arxiv.org/abs/1409.2823}. Background: 60 top-level items, compound items retained intact; source-TeX extraction, accessed 2026-07-29. The statement's quandle-homology reference slot is empty.
\item Supplied OPEN-TRIAGE summary mentions Carter--Kamada--Saito for ordinary quandle homology and scattered virtual-quandle/rack-space analogues, without complete references. It reports that its arXiv search located no dedicated paper on this specific topic. Source evidence NEEDS\_REVIEW requires a current primary-literature audit.
\item Dataset attribution: Ulam AI / UnsolvedMath Contributors, CC BY 4.0. Source rights NEEDS\_REVIEW: publicly accessible author/publisher copy, redistribution terms requiring review.
\end{itemize}
\end{document}
